\subsection{层的直像与逆像}

设 A 与 B 为拓扑空间，\(u: A \rightarrow B\) 为连续映射。

设 \(\mathcal{F} = (\mathcal{F}(\mathrm{U}), f_{\mathrm{UV}})\) 为 A 上的一个预层。我们如下定义 B 上的一个预层 \(\mathcal{F}'\)：对 B 的每个开集 U，令 \(\mathcal{F}'(\mathrm{U}) = \mathcal{F}(u^{-1}(\mathrm{U}))\)；对 B 的每对开集 (U, V) （满足 U \(\subset\) V 者），令 \(f_{\mathrm{UV}}' = f_{u^{-1}(\mathrm{U})u^{-1}(\mathrm{V})}\)。于是 \((\mathcal{F}'(\mathrm{U}), f_{\mathrm{UV}}')\) 是 B 上的一个预层。我们把它记作 \(u_{*}(\mathcal{F})\)，并称之为预层 \(\mathcal{F}\) 在映射 u 下的直像预层。

若 \((\mathrm{U}_{i})_{i\in\mathrm{I}}\) 是 B 的一族开集，则有 \(u^{-1}(\bigcup_{i\in\mathrm{I}}\mathrm{U}_{i})=\bigcup_{i\in\mathrm{I}}u^{-1}(\mathrm{U}_{i})\) 与 \(u^{-1}(\bigcap_{i\in\mathrm{I}}\mathrm{U}_{i})=\bigcap_{i\in\mathrm{I}}u^{-1}(\mathrm{U}_{i})\) （E, II, p. 25, prop. 3 与 4）。由此立即得出：若 F 具有层的性质 \((\mathrm{F}_{1})\)（相应地 \((\mathrm{F}_{2})\)）(I, p. 43)，则 \(u_{*}(\mathcal{F})\) 亦然。因此，一个层的直像是层。

设 \(\mathcal{F}_1\) 与 \(\mathcal{F}_2\) 为 A 上的预层，\(\varphi\colon\mathcal{F}_1\to\mathcal{F}_2\) 为预层的一个态射。于是存在唯一的预层态射 \(u_*\varphi\colon u_*\mathcal{F}_1\to u_*\mathcal{F}_2\)，使得对 B 的每个开集 U，映射 \((u_*\varphi)(U)\colon(u_*\mathcal{F}_1)(U)\to(u_*\mathcal{F}_2)(U)\) 就是映射 \(\varphi(u^{-1}(U))\colon\mathcal{F}_{1}(u^{-1}(U))\to\mathcal{F}_{2}(u^{-1}(U)).\) 若 \(\mathcal{F}_{3}\) 是 A 上的预层且 \(\psi\colon\mathcal{F}_{2}\to\mathcal{F}_{3}\) 是预层的态射，则有 \(u_{*}(\psi\circ\varphi)=u_{*}(\psi)\circ u_{*}(\varphi)\)。

设 C 为拓扑空间，\(v: B \to C\) 为连续映射。若 \(\mathcal{F}\) 是 A 上的预层，则预层 \(v_{*}(u_{*}(\mathcal{F}))\) 与 \((v \circ u)_{*}(\mathcal{F})\) 相同。若 \(\varphi: \mathcal{F}_{1} \to \mathcal{F}_{2}\) 是 A 上的预层态射，则有等式 \(v_{*}(u_{*}(\varphi)) = (v \circ u)_{*}(\varphi)\)。

例 1. --- 设 B 为拓扑空间，A 为 B 的一个子空间，\(\mathcal{F} = (\mathcal{F}(U), f_{\mathrm{UV}})\) 为 A 上的一个预层。用 \(i: A \to B\) 表示典范单射。于是有 \(i_*(\mathcal{F}) = (\mathcal{F}'(U), f_{\mathrm{UV}}')\)，其中对 B 的每个开集 U 有 \(\mathcal{F}'(U) = \mathcal{F}(U \cap A)\)，而对 B 的每对开集 (U, V) （满足 \(U \subset V\) 者）有 \(f_{\mathrm{UV}}' = f_{(\mathrm{U} \cap A)(\mathrm{V} \cap A)}\)。

现在设 \(\mathcal{G}\) 为 B 上的一个预层。我们定义预层 \(\mathcal{G}\) 在 \(u\) 下的逆像，并用 \(u^{*}(\mathcal{G})\) 表示层 \(\underline{\mathcal{C}}_{\mathrm{B}}(\mathrm{A};\mathrm{E}_{\mathcal{G}})\) （A 上的、以 B-空间 \(\mathrm{E}_{\mathcal{G}}\) 为值的 B-态射所成的层）（I, p. 45, 例 4）。它典范同构于 A-平展空间 \(\mathrm{A} \times_{\mathrm{B}} \mathrm{E}_{\mathcal{G}}\) 的截面在 A 上所成的层 (I, p. 9, prop. 3)。由此我们推出 (I, p. 52, 例 2) 一个典范同构 \(\varphi\)，它把与 \(u^{*}(\mathcal{G})\) 相伴的平展空间映到 A-空间 \(\mathrm{A} \times_{\mathrm{B}} \mathrm{E}_{\mathcal{G}}\) 上。此外，若 \(\mathcal{G}\) 是层，则对 A 的每一点 \(a\) 有一个典范双射 \(\psi_{a}: u^{*}(\mathcal{G})_{a} \to \mathcal{G}_{u(a)}\)：对 \(a\) 在 A 中的每个开邻域 U 与每个 B-态射 \(f: \mathrm{U} \to \mathrm{E}_{\mathcal{G}}\)，我们有

\[
\varphi ([ \mathrm{U}, f, a ]) = (a, f (a)), \quad \psi_ {a} ([ \mathrm{U}, f, a ]) = f (a).
\]

设 \(\mathcal{G}_1\) 与 \(\mathcal{G}_2\) 为 B 上的预层，\(\varphi\colon\mathcal{G}_1\to\mathcal{G}_2\) 为预层的一个态射。由 B-平展空间的态射 \(\mathrm{E}(\varphi)\colon\mathrm{E}_{\mathcal{G}_1}\to\mathrm{E}_{\mathcal{G}_2}\) 经基变换，我们导出 A-平展空间的一个态射 \(\mathrm{A}\times_{\mathrm{B}}\mathrm{E}_{\mathcal{G}_1}\to\mathrm{A}\times_{\mathrm{B}}\mathrm{E}_{\mathcal{G}_2}\)，从而导出 A 上的层的一个态射 \(u^{*}(\varphi)\colon u^{*}(\mathcal{G}_{1})\to u^{*}(\mathcal{G}_{2})\)。设 \(\mathcal{G}_3\) 为 B 上的预层，\(\psi\colon\mathcal{G}_2\to\mathcal{G}_3\) 为预层的一个态射。于是有等式 \(u^{*}(\psi\circ\varphi)=u^{*}(\psi)\circ u^{*}(\varphi)\)。

设 C 为拓扑空间，\(v\colon B \to C\) 为连续映射。若 \(\mathcal{G}\) 是 C 上的预层，则层 \(u^{*}(v^{*}(\mathcal{G}))\) 与 \((v \circ u)^{*}(\mathcal{G})\) 典范地等同 (I, p. 5)。若 \(\varphi\colon\mathcal{G}_{1} \to \mathcal{G}_{2}\) 是 C 上的预层态射，则更有 \(u^{*}(v^{*}(\varphi)) = (v \circ u)^{*}(\varphi)\)。

注. --- I, p. 53 中的典范态射 \(\sigma_{\mathcal{G}}\colon\mathcal{G}\to\widetilde{\mathcal{G}}\) 用平展空间的语言对应于 I, p. 53 命题 2 的同构 \(j_{\mathcal{G}}\)。由此推出 \(u^{*}(\sigma_{\mathcal{G}})\) 是同构。特别地，若 \(A = B\) 且 \(u = \mathrm{Id}_A\)，则预层 \(u^{*}(\mathcal{F})\) 就是层 \(\widetilde{\mathcal{F}}\)。

例 2. --- 设 B 为拓扑空间，A 为 B 的一个子空间。用 \(i\colon\mathrm{A}\to\mathrm{B}\) 表示典范单射。对 B 上的每个层 \(\mathcal{G}\)，用 \(\mathcal{G}_{\mathrm{A}}\) 表示层 \(i^{*}(\mathcal{G})\)，并称 \(\mathcal{G}_{\mathrm{A}}\) 为由层 \(\mathcal{G}\) 在 A 上诱导的层。层 \(\mathcal{G}_{\mathrm{A}}\) 等同于层 \(\underline{\mathcal{L}}(\mathrm{A};(\mathrm{E}_{\mathcal{G}})_{\mathrm{A}})\) ，即由 \(\mathrm{E}_{\mathcal{G}}\) 在 A 上诱导的 A-平展空间的截面所成的层。

设 A 是 B 的一个开子空间，\(\mathcal{G}\) 是 B 上的一个层。按定义，层 \(\mathcal{G}_{\mathrm{A}}\) 是层 \(\widetilde{\mathcal{G}}\mid \mathrm{A}\) （由 \(\widetilde{\mathcal{G}}\) 限制到开集 A 而得者）(I, p. 43)。设 \(\sigma_{\mathcal{G}}\colon \mathcal{G}\to \widetilde{\mathcal{G}}\) 为典范同构 (I, p. 55, cor. 1)。于是 \(\sigma_{\mathcal{G}}\mid \mathrm{A}\) 是从层 \(\mathcal{G}\mid \mathrm{A}\) 到层 \(\mathcal{G}_{\mathrm{A}}\) 上的一个典范同构，称为从 \(\mathcal{G}\mid \mathrm{A}\) 到 \(\mathcal{G}_{\mathrm{A}}\) 上的典范同构。

